\section{Prime dilation and nonvanishing}\label{sec:dilation}

The preceding estimates will make the rational combinations small.
We now show how the nonzero seed constant prevents them from being
zero at every large index. The coefficient of $1$ and the beta
coefficients scale differently at indices divisible by a large prime.
After one common scaling, the seed constant survives modulo that
prime, while every beta coefficient vanishes.

For a prime $p$, write
$\Zp=\{a/b\in\Q:p\nmid b\}$, and let $v_p$ be the usual
prime valuation. Recall $\tau=s-E>0$ from \eqref{eq:tau}.

\begin{proposition}\label{prop:dilation}
Fix a positive odd integer $m$. For all sufficiently large primes $p$,
\begin{align}
 p^\tau\rho_i(mp)&\in p^i\Zp,\label{eq:rho-dilation}\\
 p^\tau q_b(mp)&\equiv
       \ch(p)(m+1)^{-S}q_b(m)\pmod p.\label{eq:q-dilation}
\end{align}
In particular, the second congruence at $m=1$ transfers the seed
constant with the nonzero factor $\ch(p)2^{-S}$.
\end{proposition}

\subsection{Dilation of the rational function}

We first compare the rational functions before extracting their
coefficients. In the following calculations $p$ is large enough
that it divides none of the fixed denominators, and $p>m+1$.
Both $(mp)!$ and $(mp+1)!$ have valuation $m$. After removing
$p^m$, their units reduce to $(-1)^m m!$, by Wilson's theorem.
Fermat's theorem treats the fixed powers in $C_{mp}$. Hence
$C_{mp}$ has valuation $(S-DL)m$, and its unit reduces to
\begin{equation}\label{eq:C-unit}
 (-1)^{m(S-DL)}\frac{C_m}{(m+1)^S}.
\end{equation}

In a numerator block with fraction class $a/d$, the factors of
$P_{mp}(pz)$ divisible by $p$ are selected by
$dj+a\equiv0\pmod p$. After division by $p$, their fraction class
is $a'/d$, where $pa'\equiv a\pmod d$ and $1\le a'\le d$.
This permutes the reduced classes of each denominator and hence
permutes $\Theta$.

For example, division by $5$ interchanges the classes $1/3$ and
$2/3$: the first selected factors have
$(3+1/3)/5=2/3$ and $(1+2/3)/5=1/3$. Keeping all reduced classes
ensures that such a permutation leaves the numerator construction
unchanged.

Each block has $Lm$ selected factors. Together with the extra factor
$p(z-rm)$, their product is $p^{DLm+1}P_m(z)$. The unselected
numerator units contribute $(-1)^{DLm}$, since each complete block
of nonzero residues has product $-1$.

In the denominator, the selected indices are $h=py$, $0\le y\le m$.
They give $p^{S(m+1)}\prod_{y=0}^m(z+4y)^S$. The unselected
denominator unit is $(-1)^{Sm}$. The signs cancel:
\[
 m(S-DL)+DLm-Sm=0.
\]
The net valuation is
$(S-DL)m+(DLm+1)-S(m+1)=-s$. Thus, after removing these powers,
the resulting rational functions have the congruence
\begin{equation}\label{eq:R-dilation}
 p^s R_{mp}(pz)\equiv(m+1)^{-S}R_m(z)\pmod p.
\end{equation}
Homogeneity, \eqref{eq:homogeneity}, changes the exponent by $E$:
\begin{equation}\label{eq:T-dilation}
 p^\tau T_{mp}(pz)\equiv(m+1)^{-S}T_m(z)\pmod p.
\end{equation}
This is why the multiplier affects the prime-dilation exponent,
even though it leaves the denominator exponent and exponential
decay estimate unchanged.

\subsection{Extracting coefficients at the poles}

To deduce coefficient congruences from \eqref{eq:T-dilation}, we
must know that reduction modulo $p$ commutes with the local series
expansion. Fix a selected pole $h=py$ and put $z=-4y+u$.
The denominator factorization above shows that
\begin{equation}\label{eq:local-p-integrality}
 u^S p^\tau T_{mp}(p(u-4y))\in\Zp[[u]].
\end{equation}
After the selected powers of $p$ and the pole factor $u^S$ are
removed, all other denominator factors have unit constant term.
Their inverses therefore have coefficients in $\Zp$. The same
factorization reduces to the local expansion of the right side
of \eqref{eq:T-dilation}.

The coefficient of $u^{S-i}$ in \eqref{eq:local-p-integrality}
is $p^{\tau-i}a_{i,py}(mp)$. We may therefore extract it before
or after reduction, obtaining
\begin{equation}\label{eq:selected-coefficients}
 p^{\tau-i}a_{i,py}(mp)
       \equiv(m+1)^{-S}a_{i,y}(m)\pmod p.
\end{equation}

Now consider a pole $h=x+py$ with $1\le x<p$. Substitute
$t=pz-4x$. The extra numerator factor is a unit, the numerator
blocks have $DLm$ selected factors, and the denominator has $m$
selected factors, each of multiplicity $S$. Including $C_{mp}$,
the net valuation is zero. The same local expansion at $z=-4y$
therefore gives, for the base function,
$p^{-i}a^{\rm base}_{i,h}(mp)\in\Zp$.

Multiplication by $W_{mp}$ preserves the bound: its $k$th Taylor
coefficient is an integer and multiplies a base coefficient of order
$i+k$, with valuation at least $i+k$. Consequently
\begin{equation}\label{eq:other-coefficients}
 v_p(a_{i,h}(mp))\ge i\qquad(p\nmid h).
\end{equation}
Summing \eqref{eq:selected-coefficients} and
\eqref{eq:other-coefficients} over all poles proves
\eqref{eq:rho-dilation}.

\subsection{The surviving constant term}

Apply the constant formula \eqref{eq:q}, taking $p>4Bm$ as well.
Every denominator $j\le4bmp$ in that formula is then smaller than
$p^2$, so $v_p(j)\le1$. Multiply the formula by $p^\tau$.
At a pole with $p\nmid h$, \eqref{eq:other-coefficients} leaves
valuation at least $\tau>0$, even when $p\mid j$. At a selected
pole $h=py$, the terms with $p\nmid j$ have valuation at least
$i>0$, by \eqref{eq:selected-coefficients}.

Only the terms $h=py$, $j=pk$ survive. For these terms the upper
limit becomes $k\le4by$, and
$\ch(pk)=\ch(p)\ch(k)$. The surviving sum is therefore
$\ch(p)(m+1)^{-S}$ times the seed constant formula at index $m$.
This proves \eqref{eq:q-dilation} and completes the proof of
Proposition~\ref{prop:dilation}.
